% Compact reference appendices for equations and SI units used in the book.
\appendix
\chapter{روابط و فرمول‌های کلیدی}
\label{app:key-equations}

این پیوست رابطه‌هایی را گرد می‌آورد که در فصل سوم با آن‌ها روبه‌رو شدیم.
هدف این نیست که فرمول‌ها را دوباره از ابتدا اثبات کنیم؛ کافی است هنگام
مراجعه بدانیم هر رابطه چه چیزی را اندازه می‌گیرد و در چه چارچوبی باید
از آن استفاده کرد.

\section{شمارش حالت‌ها و نظریه‌ی اطلاعات}
اگر هر پیام از \(n\) جایگاه تشکیل شده باشد و در هر جایگاه یکی از
\(s\) نماد ممکن قرار گیرد، تعداد رشته‌های متمایز برابر است با:
\[
  M=s^n, \qquad \log_2 M=n\log_2 s
\]
خوداطلاعاتیِ رخداد \(x\)، با احتمال \(p(x)\)، به شکل زیر تعریف می‌شود:
\[
  I(x)=-\log_2 p(x)
\]
آنتروپی شانون میانگین خوداطلاعاتیِ خروجی‌های یک متغیر تصادفی \(X\) است:
\[
  H(X)=-\sum_i p(x_i)\log_2 p(x_i)
\]
اطلاعات متقابل میان ورودی \(X\) و خروجی \(Y\) نشان می‌دهد مشاهده‌ی
خروجی تا چه اندازه عدم‌قطعیت ما درباره‌ی ورودی را کاهش می‌دهد:
\[
  I(X;Y)=H(X)-H(X\mid Y)
\]
برای کانال با نویز سفید گاوسی و تحت فرض‌های مدل شانون--هارتلی، ظرفیت
نظری کانال از رابطه‌ی زیر به دست می‌آید:
\[
  C=B\log_2\left(1+\frac{S}{N}\right)
\]
در این رابطه، \(B\) پهنای باند بر حسب هرتز و \(S/N\) نسبت توان سیگنال
به توان نویز است. این رابطه را نباید بدون توجه به فرض‌های مدل به هر
نوع کانال یا مسیر ارتباطی تعمیم داد.

\section{اطلاعات و محدودیت‌های فیزیکی}
حد پایین ایده‌آلِ گرمای دفع‌شده هنگام پاک‌کردن منطقیِ یک بیت نامعلوم،
در دمای مطلق \(T\)، از اصل لاندائر به دست می‌آید:
\[
  Q_{\min}=k_{\mathrm B}T\ln 2
\]
این رابطه به عملیات منطقیِ برگشت‌ناپذیرِ پاک‌کردن مربوط است؛ نه اینکه
هر عملیات محاسباتی الزاماً همین مقدار انرژی مصرف کند.

\section{سیگنال و نمونه‌برداری}
اگر فاصله‌ی زمانی میان دو نمونه را \(T_s\) بنامیم، نرخ نمونه‌برداری
برابر است با:
\[
  f_s=\frac{1}{T_s}
\]
برای بازسازی ایده‌آلِ یک سیگنال محدودباند با بیشینه‌ی فرکانس
\(f_{\max}\)، شرط نظری نمونه‌برداری چنین است:
\[
  f_s>2f_{\max}
\]
در کاربرد واقعی، مؤلفه‌های فرکانسیِ خارج از باند باید پیش از نمونه‌برداری
تا حد لازم محدود شوند؛ در غیر این صورت هم‌پوشانی فرکانسی یا
\lr{aliasing} رخ می‌دهد.
یک مبدل \(N\)-بیتی می‌تواند \(2^N\) کد متفاوت تولید کند. تعداد کدها
به‌تنهایی دقت اندازه‌گیری را تعیین نمی‌کند؛ نویز، مرجع ولتاژ، طراحی
مبدل و کالیبراسیون نیز دخیل‌اند.

\section{روابط پایه‌ی الکتریکی}
اختلاف پتانسیل الکتریکی برابر تغییر انرژی پتانسیل به‌ازای هر واحد بار
است:
\[
  \Delta V=\frac{\Delta U}{q}
\]
جریان الکتریکی، نرخ عبور بار در زمان است:
\[
  I=\frac{\Delta Q}{\Delta t}
  \qquad\text{یا در حالت لحظه‌ای}\qquad
  i(t)=\frac{dQ}{dt}
\]
برای یک مقاومت اهمی، قانون اهم به شکل زیر نوشته می‌شود:
\[
  V=IR
\]
توان لحظه‌ای در یک عنصر الکتریکی از حاصل‌ضرب ولتاژ و جریان به دست
می‌آید:
\[
  p(t)=v(t)i(t)
\]
برای خازن ایده‌آل و سلف ایده‌آل نیز روابط پایه چنین‌اند:
\[
  Q=CV, \qquad U_C=\frac{1}{2}CV^2
\]
\[
  v_L=L\frac{di}{dt}, \qquad U_L=\frac{1}{2}LI^2
\]
در این روابط، یک نماد ممکن است در دو زمینه معنای متفاوتی داشته باشد:
\(C\) هم برای ظرفیت کانال و هم برای ظرفیت خازن به کار می‌رود؛
\(I\) نیز می‌تواند خوداطلاعاتی یا جریان الکتریکی باشد. زیرنویس‌ها و
زمینه‌ی رابطه مشخص می‌کنند کدام معنا منظور است.

\chapter{واحدهای فیزیکی و پیشوندها}
\label{app:si-units}
در فصل سوم از چند کمیت فیزیکی استفاده کردیم. جدول زیر یکاهای SI
مربوط به آن‌ها را کنار هم قرار می‌دهد. نماد یکاها در متن علمی معمولاً
لاتین نوشته می‌شود، حتی اگر نام کمیت فارسی باشد.

\begin{center}
\small
\begin{tabular}{|p{3.6cm}|p{3.5cm}|p{1.6cm}|}
\hline
\textbf{کمیت} & \textbf{یکای SI} & \textbf{نماد} \\
\hline
بار الکتریکی & کولن & \lr{C} \\
\hline
جریان الکتریکی & آمپر & \lr{A} \\
\hline
اختلاف پتانسیل & ولت & \lr{V} \\
\hline
مقاومت الکتریکی & اهم & \lr{$\Omega$} \\
\hline
توان & وات & \lr{W} \\
\hline
انرژی & ژول & \lr{J} \\
\hline
فرکانس & هرتز & \lr{Hz} \\
\hline
ظرفیت خازن & فاراد & \lr{F} \\
\hline
اندوکتانس & هِنری & \lr{H} \\
\hline
زمان & ثانیه & \lr{s} \\
\hline
دمای ترمودینامیکی & کلوین & \lr{K} \\
\hline
\end{tabular}
\end{center}

\section{پیشوندهای پرکاربرد}
پیشوندها به یکاها متصل می‌شوند تا مضرب‌ها و زیرمضرب‌های ده‌دهی را
کوتاه‌تر بنویسیم.
\begin{center}
\small
\begin{tabular}{|p{2.8cm}|p{2.4cm}|p{2.5cm}|}
\hline
\textbf{پیشوند} & \textbf{نماد} & \textbf{ضریب} \\
\hline
گیگا & \lr{G} & \(10^9\) \\
\hline
مگا & \lr{M} & \(10^6\) \\
\hline
کیلو & \lr{k} & \(10^3\) \\
\hline
میلی & \lr{m} & \(10^{-3}\) \\
\hline
میکرو & \lr{$\mu$} & \(10^{-6}\) \\
\hline
نانو & \lr{n} & \(10^{-9}\) \\
\hline
\end{tabular}
\end{center}
برای نمونه، \(1\,\mathrm{k\Omega}\) برابر با هزار اهم و
\(1\,\mathrm{mA}\) برابر با یک‌هزارم آمپر است. بزرگی و کوچکی حروف
در نماد پیشوندها اهمیت دارد؛ \lr{k} و \lr{M} دو پیشوند متفاوت‌اند.
